\documentclass[11pt]{article}
\usepackage[margin=1in]{geometry}
\usepackage[T1]{fontenc}
\usepackage[utf8]{inputenc}
\usepackage{booktabs}
\usepackage{array}
\usepackage{tabularx}
\usepackage{hyperref}
\usepackage{natbib}
\usepackage{amsmath}
\usepackage{amssymb}
\usepackage{xcolor}

\definecolor{vfblue}{HTML}{1F4E79}
\hypersetup{colorlinks=true,linkcolor=vfblue,citecolor=vfblue,urlcolor=vfblue}

\newcommand{\yes}{\checkmark}
\newcommand{\no}{--}

\title{Supplementary Materials\\[2pt]
\large VariantFlow: a selective-execution engine for efficient population\\
genomic computation on large variant datasets}
\author{Ehsan Estaji and Jian-Feng Mao}
\date{}

\begin{document}
\maketitle

\noindent This document is self-contained: each supplementary item (table or
note) begins on its own page and restates the context needed to read it
independently of the main text.

\vspace{1em}
\noindent\textbf{Contents}
\begin{itemize}\itemsep2pt
  \item Table S1 --- Cross-tool correctness validation (VariantFlow vs.\ VCFtools and scikit-allel).
  \item Table S2 --- Benchmark system configuration.
  \item Table S3 --- Timing and memory for the streaming population-genetics and Parquet/DuckDB workflows.
  \item Table S4 --- Missing-data validation against pixy.
  \item Table S5 --- Statistics coverage of VariantFlow relative to scikit-allel, grouped by computational class.
  \item Note S1 --- Missing-data-aware $\pi$ and $d_{XY}$.
  \item Note S2 --- Correctness validation methodology.
\end{itemize}

% ======================================================================
\newpage
\section*{Table S1: Cross-tool correctness validation}

VariantFlow outputs were compared against VCFtools~0.1.17 and scikit-allel~1.3.13
on the IGSR 1000 Genomes chr22 high-coverage dataset (1,073,369 biallelic
SNV records, 3,202 diploid samples). Correctness is assessed by exact
output match (integer fields) or maximum absolute difference (floating-point
fields). VariantFlow is judged correct only where the difference is zero (or
within machine epsilon for floating-point).

\begin{table}[ht]
\centering
\small
\begin{tabularx}{\linewidth}{lXccc}
\toprule
Statistic & Comparison tool & Records & Max $|\Delta|$ & Exact? \\
\midrule
Per-individual missingness (F\_MISS) & VCFtools & 3,202 & 0.0 & Yes \\
Per-site missingness (F\_MISS) & VCFtools & 1,073,369 & 0.0 & Yes \\
Allele frequency & VCFtools & 925,730 & 0.0 & Yes \\
Nucleotide diversity ($\pi$) & VCFtools & 925,730 & 0.0 & Yes \\
LD ($r^{2}$) & VCFtools & 925,730 & 0.0 & Yes \\
Weir--Cockerham $F_{ST}$ & VCFtools & 925,730 & 0.0 & Yes \\
Per-individual missingness & scikit-allel & 1,066,557 & 0.0 & Yes \\
Site-frequency spectrum (folded and unfolded) & scikit-allel & 5,000 & 0 & Yes \\
$\pi$, $d_{XY}$ pairwise counts (15\% missing) & pixy & 5,000 & 0 & Yes \\
\bottomrule
\end{tabularx}
\vspace{4pt}

\footnotesize{VCFtools rows are byte-identical after normalisation (Note~S2).
The scikit-allel missingness and site-frequency-spectrum rows are exact: maximum
absolute difference 0. For pixy, the integer pairwise counts on which the
estimator is defined are identical in every window, and the derived window
averages agree to within $5 \times 10^{-9}$. Full cross-check outputs are
recorded in \texttt{benchmark/reports/v30-correctness-crosschecks.md}. No
scikit-allel cross-check of allele frequency was performed; allele-frequency
correctness rests on the VCFtools parity check.}
\end{table}

% ======================================================================
\newpage
\section*{Table S2: Benchmark system configuration}

Benchmarks were measured on more than one host, and host choice is material:
VCFtools linkage disequilibrium takes 796~s on the Docker/Linux host and 445~s
on the Apple M5 Max, a 1.8-fold difference. The population-genetic results in
Table~S3 were measured on the M5 Max configuration below. The index-assisted
filtering and Parquet/DuckDB results were measured on a Docker/Linux host, and
the CHM13 FORMAT-filter benchmark used bcftools~1.23.1 rather than the version
listed below. Host, tool versions, and the exact commands are recorded per row
in the corresponding benchmark report. Timing is wall-clock, single-threaded;
peak memory is peak resident set size (RSS).

\begin{table}[ht]
\centering
\small
\begin{tabularx}{\linewidth}{lX}
\toprule
Parameter & Value \\
\midrule
Hardware & Apple M5 Max, 14-core CPU, 128~GB unified memory \\
OS & macOS 15.x (arm64) \\
VariantFlow & v1.5.0 (Zenodo doi:10.5281/zenodo.21198171) \\
bcftools & 1.21 (htslib 1.21) for the indexed-filter benchmarks; 1.23.1 for the CHM13 FORMAT-filter benchmark \\
VCFtools & 0.1.17 \\
scikit-allel & 1.3.13 \\
pixy & 2.2.1 \\
Input data & IGSR 1000 Genomes high-coverage (GRCh38); DDBJ CHM13 chr22 (T2T) \\
chr22 (IGSR) & 1,073,369 biallelic records, 3,202 samples \\
chr1 (IGSR) & 6,468,094 biallelic records, 3,202 samples \\
CHM13 chr22 (DDBJ) & 3,715 samples; T2T-CHM13 joint call (FORMAT-filter benchmark) \\
pixy validation set & 5,000 all-sites records, 40 samples, 2 populations, 15\% missing \\
Memory metric & Peak RSS via \texttt{resource.getrusage(RUSAGE\_CHILDREN)} \\
Timing & Wall-clock seconds, single-threaded execution \\
\bottomrule
\end{tabularx}
\end{table}

% ======================================================================
\newpage
\section*{Table S3: Timing and memory for streaming and repeated-query workflows}

Per-workflow wall-clock time and peak memory for VariantFlow and the baseline
tool, on the system described in Table~S2. Baseline is VCFtools~0.1.17 for all
population-genetics rows. Each workflow was correctness-checked against its
reference before timing (Table~S1). Timings for the FORMAT-filter and indexed
predicate workflows are summarised in Table~1 of the main text.

\begin{table}[ht]
\centering
\footnotesize
\setlength{\tabcolsep}{3pt}
\begin{tabularx}{\linewidth}{llrrrr}
\toprule
Workflow & Dataset & VariantFlow (s) & VCFtools (s) & VariantFlow (MB) & VCFtools (MB) \\
\midrule
Missingness & chr22 full (1.07M) & 23.75 & 201.67 & 8.6 & 4.1 \\
Missingness & chr1 full (6.47M) & 128.27 & 1031.13 & 8.7 & 4.1 \\
LD ($r^{2}$) & chr22 (925k) & 75.9 & 444.9 & 9.0 & 4.2 \\
Frequency & chr22 (925k) & 23.56 & 86.77 & 8.7 & 4.1 \\
Heterozygosity & chr22 (925k) & 96.09 & 140.16 & 9.3 & 4.3 \\
Hardy--Weinberg & chr22 (925k) & 85.77 & 105.55 & 9.0 & 4.7 \\
$\pi$ (per-site) & chr22 (925k) & 86.11 & 92.70 & 9.0 & 4.1 \\
Tajima's $D$ & chr22 (925k) & 100.15 & 127.05 & 27.1 & 13.8 \\
$F_{ST}$ (Weir--Cockerham) & chr22 (925k) & 84.90 & 92.14 & 9.3 & 4.7 \\
\bottomrule
\end{tabularx}
\vspace{4pt}

\footnotesize{Measured on the Apple M5 Max (Table~S2), VariantFlow~1.5.0 vs
VCFtools~0.1.17, single-threaded wall-clock; RSS = peak resident set size.
Replicates: each chromosome~22 row is the faster of two runs, LD is the median
of five isolated runs (spread 74.0--79.8~s), and the chromosome~1 row is a
single run. Speedups quoted in the main text are baseline/VariantFlow ratios.
VariantFlow
computes site and individual missingness in one pass, so the missingness baseline
sums VCFtools \texttt{--missing-site} and \texttt{--missing-indv}. $F_{ST}$ uses a
deterministic 50/50 sample split (identical for both tools), and LD is bounded
at \texttt{--max-distance 500} for both tools. ``925k'' is the
925{,}730-record biallelic-SNV chromosome~22 subset; missingness is timed on the
full biallelic chromosome. Every row was output-parity checked against VCFtools
(Table~S1).}
\end{table}

\vspace{1.5em}
\noindent\textbf{Repeated-query workflow (Parquet export + DuckDB).} The
export-once, query-many workflow was benchmarked separately because its cost
model differs from a single streaming pass: the VCF is converted to Parquet
once, after which repeated analytical queries run against the columnar file. The
measurements below use a deterministic stress VCF (DuckDB~1.5.2,
bcftools~1.21, Parquet row-group size 65\,536, five repeated queries; each time
is the wall-clock mean of three runs). The \emph{query-only} speedup is the
steady-state ratio once the Parquet file exists; the \emph{amortized} speedup
charges the full one-time export against a single query, a conservative figure
that approaches the query-only ratio as queries are repeated. Both ratios are
reproducible from the three time columns. Every query returned counts identical
to the bcftools baseline.

\begin{table}[ht]
\centering
\small
\begin{tabularx}{\linewidth}{Xrrrrrr}
\toprule
Query & Records & Export (s) & DuckDB (s) & bcftools (s) & Query-only & Amortized \\
\midrule
\texttt{QUAL>30}      & 100\,000    & 0.091 & 0.068 & 0.935  & 13.74x  & 5.87x  \\
\texttt{QUAL>30}      & 1\,000\,000 & 0.767 & 0.068 & 8.444  & 124.41x & 10.12x \\
\texttt{INFO/DP>40}   & 100\,000    & 0.085 & 0.061 & 1.605  & 26.27x  & 10.97x \\
\texttt{INFO/DP>40}   & 1\,000\,000 & 0.749 & 0.068 & 15.691 & 231.60x & 19.20x \\
\texttt{FILTER=PASS}  & 100\,000    & 0.085 & 0.066 & 0.899  & 13.58x  & 5.94x  \\
\texttt{FILTER=PASS}  & 1\,000\,000 & 0.763 & 0.070 & 9.307  & 133.66x & 11.17x \\
group-by \texttt{CHROM} & 100\,000    & 0.092 & 0.082 & 0.207  & 2.51x   & 1.19x  \\
group-by \texttt{CHROM} & 1\,000\,000 & 0.766 & 0.131 & 1.818  & 13.92x  & 2.03x  \\
\bottomrule
\end{tabularx}
\vspace{4pt}

\footnotesize{\texttt{bcftools (s)} is the repeated-scan baseline for the same
query workload. Query-only speedup $=$ bcftools\,$/$\,DuckDB; amortized speedup
$=$ bcftools\,$/$\,(export\,$+$\,DuckDB). Amortized speedup ranges from 1.19x
(grouped \texttt{CHROM} counts, 100k records) to 19.20x (\texttt{INFO/DP>40},
1M records), the 1.19--19.20x range quoted in the main text. Grouped counts
amortize weakly because a single grouped scan in bcftools is already cheap; the
workflow's advantage grows with query selectivity and repetition.}
\end{table}

% ======================================================================
\newpage
\section*{Table S4: Missing-data validation against pixy}

The pixy tool \citep{pixy2021} provides unbiased estimators for nucleotide
diversity ($\pi$) and between-population divergence ($d_{XY}$) that account
for missing genotypes by counting only non-missing pairwise comparisons at
each site, aggregated as a ratio of sums across all sites in a window
(including invariant sites). VariantFlow implements the identical estimator
in its \texttt{pixy} subcommand (Note~S1). To validate this, we generated a
deterministic all-sites VCF (5,000 sites, 40 diploid samples, two
populations) with 15\% of genotypes set to missing, and computed windowed
$\pi$ and $d_{XY}$ (1\,kb windows) with both tools.

\begin{table}[ht]
\centering
\small
\begin{tabularx}{\linewidth}{lXccc}
\toprule
Statistic & Comparison & Windows & Count match & Max $|\Delta|$ \\
\midrule
$\pi$ (per population)   & pixy 2.2.1 & 10 & Exact & $<5\times10^{-9}$ \\
$d_{XY}$ (per pop.\ pair) & pixy 2.2.1 &  5 & Exact & $<5\times10^{-9}$ \\
\bottomrule
\end{tabularx}
\vspace{4pt}

\footnotesize{Integer pairwise counts (\texttt{count\_diffs},
\texttt{count\_comparisons}) are byte-identical to pixy across every window;
the residual $|\Delta|$ in the averaged ratio reflects only the eight-decimal
print precision of VariantFlow. The comparison is performed on data
containing missing genotypes, the regime pixy was designed for.}
\end{table}

% ======================================================================
\newpage
\section*{Table S5: Statistics coverage relative to scikit-allel}

Population-genetics statistics grouped by computational class. The organizing
rule is the memory-access pattern each statistic requires. \textbf{Group A}
statistics are computable in a single streaming pass over sites or windows and
are implemented by VariantFlow's selective-execution engine; overlapping
statistics are benchmarked against scikit-allel~\citep{scikitallel} and/or
VCFtools. \textbf{Group B} statistics require the full genotype or phased
haplotype matrix to be held in memory simultaneously; these are the domain of
matrix-oriented libraries such as scikit-allel, and VariantFlow defers to them
by design rather than reimplementing them. This delineation reflects
VariantFlow's scope: it complements, and does not replace, scikit-allel.

\begin{table}[ht]
\centering
\small
\begin{tabularx}{\linewidth}{Xllc}
\toprule
Statistic & VariantFlow & scikit-allel name & Validated \\
\midrule
\multicolumn{4}{l}{\textit{Group A --- site-wise / window-wise (streaming; VariantFlow's domain)}}\\
\midrule
Allele frequency / counts & \texttt{freq} & \texttt{allele\_frequencies} & \yes\ (VCFtools, s-a) \\
Genotype missingness & \texttt{missingness} & (genotype masks) & \yes\ (VCFtools, s-a) \\
Observed/expected heterozygosity & \texttt{het} & \texttt{heterozygosity\_*} & VCFtools \\
Hardy--Weinberg equilibrium & \texttt{hardy} & \no & VCFtools \\
Nucleotide diversity $\pi$ & \texttt{pi} & \texttt{sequence\_diversity} & \yes\ (VCFtools) \\
Missing-data-aware $\pi,\,d_{XY}$ & \texttt{pixy} & (pixy) & \yes\ (pixy, exact) \\
Tajima's $D$ & \texttt{tajima-d} & \texttt{tajima\_d} & VCFtools \\
Watterson's $\theta$ & \texttt{tajima-d}$^{\dag}$ & \texttt{watterson\_theta} & partial \\
Site frequency spectrum (folded/unfolded) & \texttt{sfs} & \texttt{sfs}, \texttt{sfs\_folded} & \yes\ (s-a, exact) \\
$F_{ST}$ (Weir--Cockerham, Hudson) & \texttt{fst} & \texttt{*\_fst} & \yes\ (VCFtools) \\
Linkage disequilibrium $r^{2}$ & \texttt{ld} & \texttt{rogers\_huff\_r} & \yes\ (VCFtools) \\
\midrule
\multicolumn{4}{l}{\textit{Group B --- matrix-wide / haplotype-based (defer to scikit-allel by design)}}\\
\midrule
Joint / scaled SFS (2D) & \no & \texttt{joint\_sfs}, \texttt{sfs\_scaled} & --- \\
Selection scans (EHH, iHS, XP-EHH, nSL) & \no & \texttt{ehh}, \texttt{ihs}, \dots & --- \\
Garud's $H$, haplotype diversity & \no & \texttt{garud\_h}, \dots & --- \\
$D$-/$f$-statistics (ABBA-BABA, $f_{2/3/4}$) & \no & \texttt{patterson\_d}, \dots & --- \\
Distance, PCoA, PCA & \no & \texttt{pca}, \texttt{pcoa}, \dots & --- \\
\bottomrule
\end{tabularx}
\vspace{4pt}

\footnotesize{s-a $=$ scikit-allel. \yes\ $=$ output validated (exact or within
machine epsilon) against the named tool. A bare tool name (no \yes) means the
output was compared against that tool and agreed but is not among the quantified
rows of Table~S1; ``partial'' means only the Watterson's~$\theta$ component of
the \texttt{tajima-d} output is exercised. $^{\dag}$Watterson's $\theta$ is
reported as a component of the \texttt{tajima-d} windowed output. Group B statistics
require the whole genotype/haplotype matrix in memory and are outside
VariantFlow's streaming scope; users should use scikit-allel for these.}
\end{table}

% ======================================================================
\newpage
\section*{Note S1: Missing-data-aware $\pi$ and $d_{XY}$}

VariantFlow's \texttt{pixy} subcommand consumes an all-sites VCF and, at each
site, tallies pairwise allelic differences and comparisons over non-missing
gametes only: for a population with per-allele counts $c_a$ and
$n=\sum_a c_a$ non-missing gametes, comparisons $=\binom{n}{2}$ and
differences $=(n^2-\sum_a c_a^2)/2$; for a population pair with counts
$c^{(1)}_a,c^{(2)}_a$, comparisons $=n_1 n_2$ and differences
$=n_1 n_2-\sum_a c^{(1)}_a c^{(2)}_a$. Windowed $\pi$ and $d_{XY}$ are the
ratio of summed differences to summed comparisons over all sites in the
window. Because invariant sites contribute comparisons with zero differences
and missing genotypes are excluded per site, the estimator is unbiased under
non-uniform missingness, matching pixy~\citep{pixy2021} exactly (Table~S4)
while running as a compiled selective-execution pass.

% ======================================================================
\newpage
\section*{Note S2: Correctness validation methodology}

Each benchmark row in Table~1 of the main text undergoes correctness
validation before any speedup claim is reported. The validation process:

\begin{enumerate}
\item Run VariantFlow and the baseline tool (bcftools or VCFtools) on
  identical input with equivalent parameters.
\item Normalize outputs: strip headers, sort by genomic position, unify
  column separators.
\item Compute \texttt{diff} on normalized outputs. A row earns
  \checkmark\ only if the diff is empty (zero differing lines).
\item For floating-point statistics (e.g., $F_{\text{MISS}}$, $\pi$),
  additionally verify maximum absolute difference $<10^{-10}$.
\end{enumerate}

For the scikit-allel cross-check, we independently compute per-individual
missingness, allele frequency, and the site frequency spectrum from the
genotype matrix and compare against VariantFlow's output. This provides a
third-party validation independent of both VariantFlow and VCFtools codebases.

% Self-contained bibliography: the two works cited above are embedded directly
% so this document compiles correctly on its own (no bibtex/references.bib needed).
\begin{thebibliography}{}
\bibitem[Korunes and Samuk(2021)]{pixy2021}
Korunes, K.~L. and Samuk, K. (2021).
\newblock pixy: Unbiased estimation of nucleotide diversity and divergence in
  the presence of missing data.
\newblock \emph{Molecular Ecology Resources}, 21:1359--1368.

\bibitem[Miles et~al.(2024)]{scikitallel}
Miles, A. et~al. (2024).
\newblock scikit-allel: a {P}ython package for exploring and analysing genetic
  variation data.
\newblock \url{https://github.com/cggh/scikit-allel}.
\end{thebibliography}

\end{document}
